% !TeX program = xelatex
\documentclass[11pt,a4paper]{article}
\usepackage[margin=23mm]{geometry}
\usepackage{fontspec}
\setmainfont{Latin Modern Roman}
\usepackage{amsmath,amssymb,mathtools}
\usepackage{xcolor,enumitem,booktabs,longtable,array,xurl,hyperref}
\definecolor{muted}{HTML}{53636C}
\hypersetup{colorlinks=true,urlcolor=blue,linkcolor=blue}
\setlength{\parindent}{0pt}
\setlength{\parskip}{5pt}
\newcommand{\R}{\mathbb R}
\newcommand{\E}{\mathbb E}
\newcommand{\N}{\mathbb N}
\newcommand{\ind}{\mathbf 1}
\newcommand{\Var}{\operatorname{Var}}
\newcommand{\Cov}{\operatorname{Cov}}
\newcommand{\supp}{\operatorname{supp}}
\DeclareMathOperator*{\argmin}{arg\,min}
\DeclareMathOperator*{\argmax}{arg\,max}
\newcommand{\entrymeta}[3]{{\sffamily\small\color{muted}\textbf{\texttt{#1}}\quad|\quad #2\par #3\par}}
\newcommand{\Problemref}[1]{\texttt{#1}}
\newcommand{\JELref}[2]{\texttt{#2} (JEL \texttt{#1})}
\begin{document}
\section*{Mechanisms with resale}
\textbf{Problem ID:} \texttt{D82-78}\quad \textbf{JEL:} \texttt{D82}\par
% BEGIN ECONOMICS STATEMENT
\label{problem:OP-0292}
{\sffamily\small\color{muted}Original subject group: Mechanism design and market design\par}
\label{definitions:problem:30007574}
\textbf{Audit classification:} Published research agenda. Bergemann–Välimäki §8 discusses the interaction with resale. The optimization over initial mechanisms and disclosure at a fixed resale protocol is editorial and must not omit the continuation equilibrium.
\subsection*{Definitions and precise research target}
\textbf{Model and research target.} A seller initially owns one indivisible good. There are \(n\geq2\) buyers with private persistent values \(v_i\) drawn from a specified, common-knowledge joint distribution \(F\) on \([0,1]^n\). For each optimization, fix the resale protocol \(R\) and the remaining primitives described below. The seller commits to an initial direct mechanism \(M\) and a disclosure rule \(d\). Write \(x_i(r)\) for the probability that buyer \(i\) initially receives the object at report vector \(r\), and \(p_i(r)\) for its expected initial payment to the seller. The mechanism-induced observations, as well as disclosure signals, enter the resale game. Let \(A\) denote a candidate assessment of the complete participation, reporting and resale game, with truthful reporting and initial participation on its prescribed path. Define \(U_i^A(t,r_i;M,d,R)\) as buyer \(i\)'s interim expected total utility at true type \(t\) from participating, reporting \(r_i\), and choosing its subsequent contingent resale actions optimally against the other players' strategies in \(A\). The target is
\[
 \sup_{M,d}E_F\left[\sum_i p_i(v)\right]
 \quad\text{subject to}\quad
 U_i^A(t,t;M,d,R)\geq U_i^A(t,r_i;M,d,R)
\]
for every buyer \(i\), required type \(t\in T_i\) and permitted report \(r_i\in T_i\), together with initial participation, allocation feasibility and the declared equilibrium-selection convention for \(A\). The same assessment must be used in all incentive and participation tests for a given mechanism. Characterize how optimal disclosure and allocation depend on the fixed resale protocol and value correlation, and identify when resale weakens or strengthens initial-market revenue and efficiency under explicitly stated comparisons. The supremum concerns initial seller receipts only; it does not include any seller receipts or payments in resale.

\emph{Definitions and maintained assumptions (editorial unless a source definition is named).} This is a parameterized research agenda: the following primitives must be supplied before evaluating an instance, and are held fixed while optimizing \(M,d\). They are not additional variables in the displayed supremum. Let \(T_i=\supp(F_i)\subseteq[0,1]\) be both buyer \(i\)'s type space and its direct-report space. Supply a measurable probability kernel \(K_i(t,\cdot)\) on \(T_{-i}\), for every \(t\in T_i\), which is a version of the conditional law under \(F\). Interim constraints at all supported types use this supplied version, including at marginal-null types; \(F\) alone need not determine those values. No independence, finite-type restriction or density assumption is imposed.

Fix the participation stage after buyers learn their own types, its other observations, and a complete refusal protocol for every participation profile. The refusal protocol specifies whether the initial mechanism is cancelled or modified, which buyers can subsequently trade, payments, seller retention and all continuation opportunities. Under the candidate assessment, let \(O_i^A(t)\) be the payoff from refusing and then choosing any remaining actions optimally, with the other players' strategies fixed. Initial interim participation means
\[
 U_i^A(t,t;M,d,R)\geq O_i^A(t),\qquad i=1,\ldots,n,\quad t\in T_i.
\]
In particular, refusal does not imply a zero outside option unless the supplied refusal protocol implies it. Changes in other players' actions following an observed refusal must follow their prescribed contingent strategies, rather than a newly selected assessment.

Fix standard Borel public and private signal spaces \(S_0,S_1,\ldots,S_n\), finite initial-transfer bounds \(\ell_i\leq u_i\), and a class \(\mathcal C\) of admissible measurable joint kernels
\[
 Q(dj,dz,ds\mid r),\qquad
 j\in\{0,1,\ldots,n\},\quad
 z\in\prod_i[\ell_i,u_i],\quad
 s\in S_0\times\cdots\times S_n.
\]
Here \(j=0\) means seller retention, and \(z_i>0\) is a payment from buyer \(i\) to the seller. Write \(M\) for the ownership-and-transfer marginal of \(Q\), and \(d\) for its disclosure marginal together with the specified coupling in \(Q\). Thus \(M,d\) are admissible only with a declared joint law; their marginals alone do not determine the game. Define \(x_i(r)=Q(j=i\mid r)\) and \(p_i(r)=\int z_i\,Q(dj,dz,ds\mid r)\). Feasibility requires \(x_i(r)\geq0\) and \(\sum_i x_i(r)\leq1\). The class \(\mathcal C\), signal spaces, bounds and any restrictions on disclosure or randomization are fixed parameters, not chosen numerical conventions.

The fixed protocol and observation maps specify the order of reports, ownership and transfer realizations, signals and resale actions; who observes each report, payment, allocation and signal; remembered private histories; available offers and trading partners; bargaining power; and commitment. Allocations or observed payments can disclose reports even if the explicit signal is uninformative. The protocol also specifies whether a retaining seller can trade, its retention value, any trading costs and liquidity restrictions, and the treatment of every refusal history. All induced strategies and payoff laws must be measurable, with finite interim expected total utilities wherever a constraint is required. Buyer utility is final ownership value plus net resale receipts minus initial payments and any specified costs. Conditional expectations use \(K_i\), the joint kernel \(Q\), and the histories and strategies of the complete game.

For a hidden unilateral reporting deviation, keep the other players' contingent strategies and the candidate assessment's belief system fixed. Evaluate them at the observations actually induced by the deviation. The deviator knows its true type and remembers its report and private history, and may change all its subsequent contingent actions. No new equilibrium is selected in response to an unobserved deviation. Optimal continuation against the fixed opponents defines \(U_i^A\) as a supremum if a best response is not attained. For a truthful candidate assessment, its prescribed continuation must realize the truthful utility used in the constraint. A publicly observed change of mechanism or protocol is a different game and can have a different selected assessment.

An assessment must specify strategies and beliefs at all relevant information sets, with Bayesian updating wherever applicable and the declared off-path consistency requirements. A finite bargaining horizon does not make this complete game finite when types, signals or actions are infinite. The extension of sequential equilibrium to the admitted spaces, its consistency definition, and the selection rule over mechanism-dependent complete assessments remain to be specified. Once supplied, they must state how mechanisms without an appropriate assessment or measurable selection are treated; existence or attainment is not presumed. Similarly, revenue and efficiency comparisons require a named no-resale benchmark, a statement of whether mechanisms are held fixed or reoptimized, and a final-ownership welfare functional including any seller retention value and trading costs. These comparison choices are not determined by the displayed initial-receipts objective.
\subsection*{Variants and assumption changes}
\begin{enumerate}
\item \textbf{Posted-price resale (editorial).} Fix a designated bargaining order in which the current owner makes one take-it-or-leave-it offer, with bounded price grid and public versus private initial reports. Characterize optimal initial mechanism and disclosure under independent finite types.
\item \textbf{Correlated resale (editorial).} Admit correlated finite types and private disclosure signals; compare optimistic versus pessimistic selection among continuation equilibria. Any revenue/efficiency conclusion must identify which selection convention it uses.
\end{enumerate}
\subsection*{References and source audit}
\begin{enumerate}
\item §8 printed pp.50–51 in June2018 revision. Supports the stated research area; exact displayed specification is editorial. DOI publication is2019; inspected revision is June2018. \url{https://cowles.yale.edu/sites/default/files/2022-09/d2102-r.pdf}
\end{enumerate}
\begin{enumerate}\raggedright
\item\hypertarget{definitions:source:30007574:1}{} Dirk Bergemann; Juuso Välimäki. 2019. \emph{Dynamic Mechanism Design: An Introduction}. Published JEL57(2),235–274 (2019); inspected June19,2018 Cowles revision §8, printed pp.49–51.
\url{https://cowles.yale.edu/sites/default/files/2022-09/d2102-r.pdf}.
DOI: \url{https://doi.org/10.1257/jel.20180892}.
\end{enumerate}

\subsection*{Common conventions: Source mathematical conventions}

These conventions apply to each entry unless explicitly replaced.

\textbf{Models and identification.} A primitive model \(m\) specifies all probability laws, feasible choices, information, equilibrium and measurement rules used by its target. The admissible class \(\mathcal M\) is fixed independently of outcomes. If \(P_O(m)\) is its observable probability law and \(\tau(m)\) the target, its identified set at observable law \(P\) is
\[
 \mathcal I_\tau(P)=\{\tau(m):m\in\mathcal M,\ P_O(m)=P\}.
\]
Point identification means that this set is a singleton. An empty set signals incompatibility of data and assumptions. Coordinate bounds need not describe the joint set. Bounds are sharp only if every retained target is attainable or arbitrarily approachable by compatible models. Finite bounds require explicit restrictions on unobserved quantities; unrestricted confounding, hidden wealth, extrapolation or arbitrary equilibrium selection can yield uninformative sets.

\textbf{Causal interventions.} A potential outcome \(Y_i(p)\) is the outcome under a complete policy regime \(p\), including timing, financing, information and market spillovers. The target population and units are fixed before treatment. With interference, use the complete assignment vector or a justified exposure mapping; individual no-interference notation alone cannot identify an economy-wide intervention. A difference of mean potential outcomes does not require the cross-world joint distribution. Conditioning on post-treatment migration, survival, enrollment or employment changes the estimand and needs separate justification.

\textbf{Statistical statements.} An estimator, confidence set, forecast or test specifies a sampling law, model class and loss/coverage criterion. Uniform \((1-\alpha)\)-coverage means \(\inf_{m\in\mathcal M}\Pr_m\{\tau(m)\in C_n\}\geq1-\alpha\), or an explicitly asymptotic version. The whole parameter space trivially has coverage, so informativeness requires a stated width, risk or power objective. A forecasting improvement is relative to a named benchmark, information set and evaluation population.

\textbf{Equilibrium and optimization.} An equilibrium comprises feasible plans, optimality, beliefs where needed, budget constraints and all market/resource clearing. The primitives or a stated benchmark define each correspondence \(\mathcal E(p;m)\). Nonemptiness is assumed or proved, never inferred just from compact policy sets. A selected-equilibrium target uses a declared selection rule; a robust target quantifies over all permitted equilibria. Use suprema or infima unless attainment is established. An \(\varepsilon\)-optimal policy is feasible and within \(\varepsilon\) of the finite optimum value. Report an empty feasible set separately.

\textbf{Welfare and accounting.} Normative utility, interpersonal normalization, social weights, numeraire, discounting and the use of public receipts are primitives. Behavioral utility need not coincide with the welfare criterion. Household welfare includes ownership income and rents once; adding profits again to compensating variation generally double-counts them. A monetary equivalent is defined by an explicit transfer and price convention, with monotonicity and range conditions. Average effects, mechanism contributions, transfers and welfare losses are different targets. Infinite sums require convergence and dynamic feasible plans require terminal or no-Ponzi conditions.

\textbf{Source versions.} A working-paper date, revision date and journal publication year are distinguished. A DOI for a journal version is not automatically the DOI of an earlier manuscript. Locators refer to the stated version and printed pages where specified. A metadata-only check does not verify a claim about a full-text conclusion. Every entry's source subsection narrows the provenance claim accordingly.
% END ECONOMICS STATEMENT
\end{document}
